\subsection{局部紧齐性空间}

命题 13. 设 G 是局部紧群，H 是 G 的闭子群。则齐性空间 G/H 是局部紧的且是仿紧的。

由于 G/H 是豪斯多夫的（§ 2，第 5 节，命题 13），把第 5 节命题 9 应用于右作用于 G 的 H，可知它是局部紧的。于是只需再证 G/H 是仿紧的。设 V 是 G 中 e 的对称紧邻域，并设 \(G_{0}=V^{\infty}\) 是由 V 生成的 G 的子群。\(G_{0}\) 是开的（§ 2，第 1 节，命题 4 的推论）且连续作用于 G/H（§ 2，第 5 节，命题 12）。如果我们能证明每个轨道 \(G_{0}.z\)（ \(z\in G/H\)）都是 G/H 的开子集且是紧集的可数并，那么就可推出 G/H 是诸不同轨道 \(G_{0}.z\) 的拓扑和，从而是仿紧的（第一章，§ 9，第 10 节，定理 5）。\(\mathbf{G}_0.z\) 在 G/H 中是开的，这由 \(\mathbf{G}_0\) 在 G 中是开的以及等价关系 \(x^{-1}y\in H\) 是开的（§ 2，第 4 节，引理 2）这两个事实得出。另一方面，\(\mathbf{G}_0.z\) 是诸集合 \(\mathbf{V}^n.z\)（ \(n\geqslant 1\)）的并，且由于 \(\mathbf{V}^n\) 在 G 中紧而 G/H 是豪斯多夫的，\(\mathbf{V}^n.z\) 是紧的。证毕。

命题 14. 在局部紧群 G 中，单位元连通分支 C 是 G 的诸开子群的交。

C 是 G 的闭正规子群（§ 2，第 2 节，命题 7），从而 G/C 是局部紧的（命题 13）且全不连通的（第一章，§ 11，第 5 节，命题 9）。由于 G 到 G/C 上的典范映射下 G/C 的开子群的逆像是包含 C 的 G 的开子群，可见只需对群 G/C 证明本命题。换句话说，我们可以假设 G 是全不连通的。此时我们知道（第二章，§ 4，第 4 节，命题 6 的推论），e 的每个紧邻域都包含 e 的一个既开又闭的邻域 U。由于 U 是紧的且 \(B = \complement U\) 是闭的，存在 e 的对称开邻域 W 使得 W ⊂ U 且 UW ∩ BW = ∅（§ 3，第 1 节以及第二章，§ 4，第 3 节，命题 4），更有 UW ⊂ U。对 n 用归纳法推出，\(W^n \subset U\) 对每个整数 \(n > 0\) 成立。于是由 W 生成的子群 \(W^\infty = \bigcup_{n > 0} W^n\) 包含在 U 中；但 \(W^\infty\) 在 G 中是开的（§ 2，第 1 节，命题 4 的推论）。证毕。

我们同时也证明了：

推论 1. 若 G 是全不连通的局部紧群，则 G 中 e 的每个邻域都包含 G 的一个开子群。

推论 2. 局部紧群若由单位元的每个邻域生成，则它是连通的。

推论 3. 设 G 是局部紧群，H 是 G 的闭子群，\(\varphi\) 是 G 到 G/H 上的典范映射。则 G/H 的连通分支是 G 的连通分支在 \(\varphi\) 下的像的闭包。

设 C 是 G 的单位元连通分支。G 的连通分支是诸集合 sC，其中 \(s \in G\)（§ 2，第 2 节，命题 7）；\(\varphi(sC)\) 显然是连通的，从而 \(\overline{\varphi(sC)}\) 也是连通的（第一章，§ 11，第 1 节，命题 1）。但是 \(\varphi(sC) = \varphi(sCH)\)，且由于 sCH 关于由 H 定义的等价关系是饱和的，又由于这个等价关系是开的

（§ 2，第 4 节，引理 2），我们有 \(\overline{\varphi(s\mathbf{CH})} = \varphi(\overline{s\mathbf{CH}}) = \varphi(s.\overline{\mathbf{CH}})\)（第一章，§ 5，第 3 节，命题 7）。

令 L = CH。L 是 G 的闭子群且包含 C 与 H；因此为证诸集合 \(\varphi(s.L) = s.\varphi(L)\) 是 G/H 的连通分支，只需证明 G/H 按以诸集合 \(s.\varphi(L)\) 为类的等价关系所得的商空间是全不连通的。而这个商空间同胚于齐性空间 G/L（第一章，§ 3，第 4 节，命题 7）；于是问题归结为证明当 \(C \subset H\) 时 G/H 是全不连通的。由于 G/H 可与 \((\mathrm{G/C})/(\mathrm{H/C})\) 等同（§ 2，第 7 节，命题 22），我们甚至可以假设 G 本身是全不连通的。G/H 中 \(\varphi(e)\) 的每个邻域都包含形如 \(\varphi(V)\) 的邻域，其中 V 是 G 中 e 的邻域，从而（由推论 1）包含形如 \(\varphi(K)\) 的邻域，其中 K 是 G 的紧开子群。于是 \(\varphi(K)\) 在 G/H 中既开又闭，这表明 G/H 中 \(\varphi(e)\) 的连通分支仅由 \(\varphi(e)\) 一点组成。通过平移，对 G/H 的每一点的连通分支同样成立，推论得证。

